% Generated draft: requires mathematical and source review.
\documentclass[11pt]{article}
\usepackage[margin=25mm]{geometry}
\usepackage{fontspec}
\setmainfont{Latin Modern Roman}
\usepackage{amsmath,amssymb,mathtools,unicode-math}
\setmathfont{Latin Modern Math}
\usepackage{enumitem,needspace,xurl,hyperref}
\setlength{\parindent}{0pt}
\setlength{\parskip}{5pt}
\setlength{\emergencystretch}{3em}
\newcommand{\N}{\mathbb{N}}
\newcommand{\Z}{\mathbb{Z}}
\newcommand{\Q}{\mathbb{Q}}
\newcommand{\R}{\mathbb{R}}
\newcommand{\C}{\mathbb{C}}
\newcommand{\F}{\mathbb{F}}
\newcommand{\A}{\mathbb{A}}
\newcommand{\PP}{\mathbb{P}}
\newcommand{\E}{\mathbb{E}}
\newcommand{\eps}{\varepsilon}
\newcommand{\dd}{\,\mathrm{d}}
\begin{document}
\section*{UnsolvedMath ID: 800949}
\textbf{Source problem code:} catalogue entry
\subsection{Definitions and mathematical statement}
Let \(g\) be a complete Riemannian metric on \(\R^2\), and let \(\nabla^g\) denote its Levi--Civita connection. Interpret a non-trivial geodesic as the image of a complete nonconstant geodesic, with parametrization and orientation disregarded. Equivalently, define
\[
\mathcal{G}_g=\left\{\gamma(\R):\gamma:\R\longrightarrow\R^2,\quad \nabla^g_{\dot\gamma}\dot\gamma=0,\quad g(\dot\gamma,\dot\gamma)=1\right\}.
\]
Completeness ensures that every unit-speed geodesic extends over all of \(\R\). The use of complete geodesic images makes the uniqueness assertion independent of reparametrization; the supplied source does not explicitly specify this convention.

Assume the following version of Euclid's fifth postulate:
\[
\forall L\in\mathcal{G}_g\ \forall p\in\R^2\setminus L,\qquad
\#\left\{L'\in\mathcal{G}_g:p\in L'\ \text{and}\ L'\cap L=\varnothing\right\}=1.
\]
Must \(g\) be flat, that is, must its Gaussian curvature satisfy
\[
K_g(x)=0\qquad\text{for every }x\in\R^2?
\]
The supplied source additionally states that the parallel-postulate assumption implies the absence of conjugate points. This is retained as a source assertion, rather than imposed as a separate hypothesis or independently verified here.
\subsection{Short English statement}
On the plane with a complete Riemannian metric, suppose that through every point outside any complete geodesic there is exactly one complete geodesic disjoint from it, counting geodesics by their images. Must the metric have zero Gaussian curvature everywhere? The supplied registry labels this question open; this draft does not independently verify that status.
\subsection{Sources}
\begin{itemize}
\item Local UnsolvedMath registry, permanent ID \(800949\), source problem code \texttt{catalogue entry}; supplied title: \emph{Open Problems and Questions about Geodesics --- Question 29 ([Bur-Kni1991])}. The record belongs to the AMR Open Problem Lists collection and supplies the status \texttt{open}.
\item Burns and Matveev, \emph{Open Problems and Questions about Geodesics}, listed in the supplied record as a 2019 source, Question 29. Supplied source URL: \url{https://arxiv.org/abs/1308.5417}. The record attributes the question to Burns and Matveev and gives proposed year 2019; these are supplied attribution fields.
\item The background gives the attribution/context label \texttt{[Bur-Kni1991]}, without a full bibliographic entry. Its authors, title, and publication details cannot be established from the supplied record.
\item The background records access on 2026-10-16 and extraction by \texttt{arxiv\_latex}. These are supplied provenance notes, not a claim that this draft independently accessed or verified the source. The registry records creation and update on 2026-10-16 and upstream revision \texttt{d371ee917ac9647fbb06cd6fc529deec6f3dda82}.
\end{itemize}
\end{document}
